% =====================================================================
%  Chapter 8: Dynamical fermions
%  Gattringer & Lang, "Quantum Chromodynamics on the Lattice" (LNP 788)
%  Reconstructed from extract/ch08.txt; equations auto-numbered.
% =====================================================================

\chapter{Dynamical fermions}

QCD is the quantum field theory of gluons and quarks. In previous chapters
we have found that for gluons the inclusion of the full dynamical behavior
in Monte Carlo simulations is straightforward, but not so for quarks. In this
chapter we discuss the still missing feature: dynamical quarks.
A calculation with dynamical fermions is a much more challenging enterprise
than a quenched calculation, and new algorithmic ideas are needed for the
inclusion of the fermion determinant. In this chapter we present some of these
algorithms and techniques. Subsequently we address the coupling-mass phase
diagram and discuss some results of dynamical QCD calculations.

\section{The many faces of the fermion determinant}

In the path integral, whenever the fermions occur bilinearly in the action, the
Grassmann integral can be evaluated in closed form and leads to the fermion
determinant as a factor. As an example we quote the expression for the
two-point function of a flavor triplet meson from Sect.~6.1.5,
\begin{equation}
\begin{split}
\ev{O_T(n)\,O_T(m)} &= -\frac{1}{Z}\int D[U]\,e^{-S_G[U]}\det[D_u]\det[D_d]\\
&\qquad\times \tr\!\big[\Gamma D_u^{-1}(n|m)\,\Gamma D_d^{-1}(m|n)\big],\\
Z &= \int D[U]\,e^{-S_G[U]}\det[D_u]\det[D_d].
\end{split}
\end{equation}
The meson interpolators are converted into products of quark propagators and
each fermion flavor accounts for one fermion determinant factor. In our example
we take into account the two lightest quarks $u$ and $d$. In the quenched
approximation, discussed in Sect.~6.1.6, the determinants were put to unity by
hand. This amounts to neglecting vacuum loops of quarks (see Fig.~6.2).
Here we now take into account the effect of the fermions. We stress that the
fermion determinant is a functional of the gauge field and has to be computed
anew for every gauge configuration. Even for lattices of moderate size the
Dirac operator is a huge matrix with $N = 12|\Lambda|$ rows and columns, where
12 is the product of color and Dirac entries and $|\Lambda|$ the total number
of lattice points. This number easily becomes larger than a million. The
computational cost of calculating even only one such determinant in closed
form, which formally has $N!$ contributing terms, is prohibitively high.

\subsection{The fermion determinant as observable}

Let us forget for a moment the technical problem of high computational cost.
For small lattice size and a few configurations a direct determination is
certainly possible. In this case one can try to treat the determinant as part
of an observable and rewrite (8.1) as
\begin{equation}
\ev{O(U,\psi,\bar{\psi})} = \frac{\ev{\det[D_u]\det[D_d]\,O[U]}_G}{\ev{\det[D_u]\det[D_d]}_G}.
\end{equation}
The expectation value $\ev{\cdots}_G$ is computed with the path integral for
pure gauge theory,
\begin{equation}
\ev{A}_G = \frac{1}{Z_G}\int D[U]\,e^{-S_G[U]}A[U], \qquad
Z_G = \int D[U]\,e^{-S_G[U]}.
\end{equation}
Although the idea of treating the determinant as an observable is appealing, it
has a serious flaw: Depending on the gauge configuration the determinant may
have widely different values, typically covering several orders of magnitude.
In the sum over all configurations this leads to large fluctuations around the
mean value and thus to an intrinsic instability. As a matter of fact, the
distribution of the gauge configurations according to (8.3) is very different
from that obtained when also the determinant is taken into account in their
Monte Carlo generation.
For that reason, treating the determinant as an observable is numerically
justified only for extremely large statistics. Only in lower dimensions (e.g.,
for the 2D Schwinger model) such an approach has led to acceptable results.

\subsection{The fermion determinant as a weight factor}

In order to obtain properly distributed gauge configurations, one tries to
include the determinant as a probability weight factor when generating the
Markov chain of gauge configurations. Thus the gauge fields $U$ are distributed
according to the joint distribution (for two dynamical quark flavors)
\begin{equation}
\frac{1}{Z}\,e^{-S_G[U]}\det[D_u]\det[D_d].
\end{equation}
Once the gauge configurations have been generated according to this
distribution, the observables are computed on these gauge configurations as
discussed in Chaps.~4 and 6.

There is, however, a potential problem: If we want to interpret the contribution
of the fermion determinants as a factor in a probability weight it must be real
and nonnegative. In most cases one can use the $\gamma_5$-hermiticity of $D$
(see (5.76)), i.e., $\gamma_5 D \gamma_5 = D^\dagger$, to show that the
determinant is real:
\begin{equation}
\det[D]^* = \det[D^\dagger] = \det[\gamma_5 D\gamma_5] = \det[D].
\end{equation}
However, the determinant could still be real but negative. A possible
physically justified (see (7.21)) solution is to assume that $u$ and $d$ quarks
are mass degenerate and thus $D_u = D_d \equiv D$. More generally, for an even
number of mass-degenerate quarks the fermion determinant is raised to an even
power and the combined weight factor is nonnegative. No principal obstacle for
an interpretation as probability weight is encountered then. For the case of
two degenerate flavors we can use (8.5) and write
\begin{equation}
0 \le \det[D]\det[D] = \det[D]\det[D^\dagger] = \det[DD^\dagger].
\end{equation}
In the last step we have rewritten the product of the two determinants for the
two flavors into a single determinant of the hermitian matrix $DD^\dagger$, a
form which will be useful later.
If in addition to $\gamma_5$-hermiticity the Dirac operator is also chiral (for
vanishing quark mass), it obeys the Ginsparg--Wilson equation in the form
(7.48). The spectrum for the massless case is then restricted to the
Ginsparg--Wilson circle shown in Fig.~7.1. Since the eigenvalues come in
complex conjugate pairs (cf.\ (7.45)), the determinant is real and nonnegative.
If one introduces a mass term, the determinant is even strictly positive. Thus
for a chiral $\gamma_5$-hermitian Dirac operator also odd powers of the fermion
determinant give rise to a proper probability weight.
In special but important cases, e.g., when a chemical potential or a
$\theta$-angle are introduced, $\gamma_5$-hermiticity does not hold and the
determinant may be even complex. These cases need to be treated with different
methods (see Chap.~12).

\subsection{Pseudofermions}

The central idea in introducing dynamical fermions in the Monte Carlo sampling
of gauge fields is based on the analogy between fermionic and bosonic Gaussian
integrals. Consider the vector $\varphi = \varphi_R + i\varphi_I$ of $N$
complex variables and a matrix $A$ with eigenvalues $\lambda$ which all have a
positive real part $\Re\lambda > 0$. Then a Gaussian integral equivalent to
(5.32) may be evaluated giving the generating functional
\begin{equation}
\begin{split}
W[\chi,\chi^\dagger] &\equiv \int_{\mathbb{R}^{2N}}\Big(\prod_{i=1}^{N}d\varphi_{R,i}\,d\varphi_{I,i}\Big)
\exp\Big(-\sum_{i,j=1}^{N}\varphi_i^\dagger A_{ij}\varphi_j\\
&\qquad\qquad + \sum_{i=1}^{N}\varphi_i^\dagger\chi_i + \sum_{i=1}^{N}\chi_i^\dagger\varphi_i\Big)\\
&= \frac{\pi^N}{|\det[A]|}\exp\Big(\sum_{i,j=1}^{N}\chi_i^\dagger(A^{-1})_{ij}\chi_j\Big).
\end{split}
\end{equation}
Up to an irrelevant factor $\pi^N$, the result agrees with the corresponding
Grassmann integral, except that the determinant is now in the denominator.
For simplicity we again introduce the abbreviation
$\int D[\varphi_R]D[\varphi_I] \equiv \prod_i d\varphi_{R,i}\,d\varphi_{I,i}$.
This correspondence has led to the introduction of the so-called
pseudofermions. These are bosons with the same number of degrees of freedom as
the fermionic variables. One has used them in two different ways:
\begin{itemize}
  \item With $A \equiv (DD^\dagger)^{-1}$ one may express the fermion
        determinant as Gaussian integral of a bosonic field. Used in that
        context we call the variables pseudofermion fields (see below).
  \item With $A \equiv (DD^\dagger)$ one may utilize the Gaussian integral in
        order to estimate matrix elements like $D^{-1}(n|m)$ or sums like
        $\tr[D^{-1}]$. This Monte Carlo sum introduces statistical noise and
        therefore we call these variables noisy pseudofermions used to
        construct noisy pseudofermion estimators. These techniques will be
        discussed in Sect.~8.4.
\end{itemize}

For a $D$ with nonvanishing eigenvalues we have
$\det[D] = 1/\det[D^{-1}]$. Thus one may further utilize the boson-fermion
analogy. In [1] the pseudofermion fields were introduced by observing that from
(8.7) one obtains
\begin{equation}
\det[DD^\dagger] = \pi^{-N}\int D[\varphi_R]D[\varphi_I]\,e^{-\varphi^\dagger(DD^\dagger)^{-1}\varphi},
\end{equation}
where we suppress color, Dirac, flavor, and lattice indices and use
vector/matrix notation for $\varphi$ and $D$. One uses this relation to replace
the integral over the fermionic Grassmann variables for two mass-degenerate
flavors by an integral over bosonic variables,
\begin{equation}
\int D[\psi]D[\bar{\psi}]\,e^{-\bar{\psi}_u D\psi_u - \bar{\psi}_d D\psi_d}
= \pi^{-N}\int D[\varphi_R]D[\varphi_I]\,e^{-\varphi^\dagger(DD^\dagger)^{-1}\varphi}.
\end{equation}
In this formulation the number of fermions of a given mass has to be even in
order to guarantee positivity which is necessary for the convergence of the
Gaussian integrals. The effective bosonic interaction contains the inverse
matrix $(DD^\dagger)^{-1}$ and therefore is highly nonlocal.

\subsection{Effective fermion action}

A useful way of thinking about the fermion determinant is to interpret it as an
additional contribution to the gauge action, the so-called effective fermion
action. With (A.54) we write the determinant in the form of an exponential and
obtain
\begin{equation}
\det[D] = \exp(\tr[\ln(D)]) \equiv \exp(-S_F^{\mathrm{eff}}), \qquad
S_F^{\mathrm{eff}}[U] = -\tr[\ln(D)].
\end{equation}
Here we have assumed that $\det[D]$ is real and positive, i.e., $D$ is either a
chiral $\gamma_5$-hermitian matrix plus a mass term, or we work with two
flavors and $D$ is replaced by $DD^\dagger$ according to (8.6).

In contrast to the local gauge action, the effective fermion action is a very
nonlocal quantity and connects essentially all gauge variables of the system
with each other. Here we use the terms nonlocal and local to distinguish
whether coupling is to all variables of the system or only to field variables
in the neighborhood. The latter is often called ultralocal to distinguish it
from being local only in the continuum limit, for which it is sufficient to
show that the coupling to distant field variables decreases exponentially with
distance (compare the discussion of the locality properties for the Dirac
operator in Sect.~7.4.2).

Generating gauge configurations distributed according to the combined
Boltzmann weight factor
\begin{equation}
\exp(-S[U]) \qquad \text{with} \qquad S[U] = S_G[U] + S_F^{\mathrm{eff}}[U]
\end{equation}
is a possible way to include dynamical fermions. For generating sequences of
gauge configurations according to the distribution (8.11) we need the change of
the total action
\begin{equation}
\exp(-S[U'] + S[U]) = \exp(-S_G[U'] + S_G[U])\,
  \exp(-S_F^{\mathrm{eff}}[U'] + S_F^{\mathrm{eff}}[U])
\end{equation}
for the Metropolis step, where we decide whether to accept the candidate
configuration $U'$ or reject it. Computing the factor from the gauge action is
simple (compare Sect.~4.1). The second term, coming from the effective fermion
action, requires new techniques.

\subsection{First steps toward updating with fermions}

Before we describe modern fermion algorithms in more detail, let us briefly
discuss some general ideas for the update. Different methods differ in the way
the action change is determined and in the dependence of the computer time on
the volume of the lattice and the fermion mass:
\begin{itemize}
  \item Global changes of the gauge field, accepting with the determinant as
        weight factor: One estimates the ratio
        \begin{equation}
        \exp(-\Delta S_F^{\mathrm{eff}})
        \equiv \exp(-S_F^{\mathrm{eff}}[U'] + S_F^{\mathrm{eff}}[U])
        = \frac{\det[D[U']]}{\det[D[U]]}
        \end{equation}
        by stochastic methods as discussed in Sect.~8.4.
  \item Small steps, estimating $\Delta S_F^{\mathrm{eff}}$ by an approximation
        essentially linear in the change of the gauge variables. Due to this
        approximation errors may accumulate and one has to rely on corrective
        measures.
\end{itemize}

Both paths have been pursued.

Updating, e.g., a single gauge link variable in a Metropolis step requires the
determination of the action change. For this one has to compute $D^{-1}$ acting
on a vector, usually employing a conjugate gradient algorithm (see
Sect.~8.2.3). Assuming that the necessary number of iterations depends only
weakly on the volume, one such inversion needs computer time proportional to
the lattice volume. The acceptance may be reasonable, but this computation has
to be repeated for each change of a gauge variable, and therefore the full
algorithmic cost grows with $V^2$.
On the other hand, considering a global update, where all link variables are
altered (or a partial-global update, cf.\ [2]), the change of the action will
be typically proportional to the volume and the acceptance rate drops
exponentially with the volume. Thus one has to decrease the step size inversely
proportional to the volume, which again leads to a computational effort growing
with the volume squared.
However, there are algorithms with much better performance, based on the
pseudofermion field concept. The idea is to choose changes of the gauge field
not completely randomly but along directions determined by the action. This
introduces a bias, but there are methods to effectively undo this bias in the
updating such that the detailed balance is respected. The so-called hybrid
Monte Carlo (HMC) algorithm is such a method and will be discussed next.

\section{Hybrid Monte Carlo}

The updating of the gauge field variables in a distribution
$P_S(U) \propto \exp(-S)$ defined through the gauge field action and the
effective fermionic action has two parts. First one has to find a reasonable
candidate for a change of the gauge variables. This introduces an a priori
selection probability factor $T_0(U'|U)$ (compare Sect.~4.1.3). In a second
step one has to decide whether to accept or reject the new configuration,
according to an acceptance probability $T_A(U'|U)$. Together the two steps
provide the overall transition probability
\begin{equation}
T(U'|U) = T_A(U'|U)\,T_0(U'|U).
\end{equation}
The detailed balance condition (4.15) for $T(U'|U)$ may be obeyed with a
Metropolis accept--reject probability following (4.16).

If one considers only the gauge field action both steps are simple since the
gauge action is ultralocal (cf.\ Chap.~4). Thus computing the change of the
gauge action is cheap and one makes small changes of variables, maybe even just
for one link, and then decides whether to accept or reject the change.
Applying this to the full action including fermions is computationally
expensive since the determinant leads to a nonlocal $S^{\mathrm{eff}}$ and thus
computing the change of the action involves all link variables, even if only a
single gauge link is altered. One therefore attempts to update many variables
in one step.

Doing this in a naive manner typically leads to large changes of the action and
to extremely small acceptance probability and thus is not efficient. Ideally,
the acceptance probability should be large and only weakly volume dependent. At
the same time the autocorrelation between subsequent configurations should be
as small as possible.
There are various strategies to find a new configuration with a reasonable
chance to be accepted. We now discuss the differential equation approach.

\subsection{Molecular dynamics leapfrog evolution}

Let us first look at the basic concept: A Hamiltonian process developing in
computer time (Markov time) $\tau$. In order to focus on the central idea we
discuss the method first for a real scalar field $Q$. For simplicity of
presentation we suppress the lattice index for this field and employ
vector/matrix notation. Using the gauge variables leads to technical
complications which will be addressed later. Conjugate to the bosonic field
variables $Q$ we introduce real momenta $P$ and write the formula for the
vacuum expectation value of some observable $O$ as
\begin{equation}
\ev{O}_Q = \frac{\int D[Q]\,e^{-S[Q]}\,O[Q]}{\int D[Q]\,e^{-S[Q]}}
= \frac{\int D[Q]D[P]\,e^{-\half P^2 - S[Q]}\,O[Q]}
        {\int D[Q]D[P]\,e^{-\half P^2 - S[Q]}}
= \ev{O}_{P,Q}.
\end{equation}
To give an example for our use of vector/matrix notation, we remark that $P^2$
stands for
\begin{equation}
P^2 = \sum_{n\in\Lambda}P(n)^2.
\end{equation}
The second form in (8.15) is equivalent since in the expectation value the
Gaussian integrals for $P$ cancel. In the large volume limit this new form
represents a microcanonical ensemble of a classical (nonrelativistic) system
with Hamiltonian
\begin{equation}
H[Q,P] \equiv \half P^2 + S[Q].
\end{equation}
The classical equations of motion (in computer time) for this system are
\begin{equation}
\dot{P} = -\frac{\partial H}{\partial Q} = -\frac{\partial S}{\partial Q},
\qquad \dot{Q} = \frac{\partial H}{\partial P} = P.
\end{equation}
The notation used for the derivatives is a shorthand notation for the
derivative with respect to the field value at a single lattice site in
accordance with our vector/matrix convention.

The Eqs.~(8.18) are called molecular dynamics equations since they determine
the time evolution of a classical system of particles. The Hamiltonian is a
constant of motion, and the path of the configurations $(P,Q)$ lies on a
hypersurface of constant energy in phase space and thus would be always
accepted if the evolution of (8.18) could be done exactly.
The Eqs.~(8.18) can be evolved numerically, and according to (8.15) we may
extract the requested expectation values $\ev{O}_Q$ from the $Q$--$P$
ensemble. We stress that one also has to show that the evolution gives rise to
an update which is ergodic for the fields of interest, i.e., the fields $Q$ in
this case.
Using the microcanonical evolution for pure gauge theories has been suggested
in [3], for fermions in [4]. Such molecular dynamics evolution equations are
also closely related to stochastic differential equations like the Langevin
equation, where the classical evolution is disturbed by a noise term. The
random noise gives rise to the quantum fluctuations of the quantum field theory
[5, 6]. Changing the conjugate momenta randomly, according to the Gaussian
distribution, and pursuing the microcanonical path in between these changes led
to the so-called hybrid algorithms [7, 8].
The numerical implementation of (8.18) introduces a discrete step size
$\varepsilon \equiv \Delta\tau$ and numerical errors are unavoidable. A simple
linear evolution scheme introduces errors $\order{\varepsilon^2}$. Either one
finds a way to extrapolate the results to vanishing step size or one introduces
a corrective step. This second idea leads to the hybrid Monte Carlo algorithm.
We call a sequence of small steps following the approximate molecular dynamics
evolution a trajectory. For the HMC method one first constructs a new
configuration by evolution along such a trajectory and then decides in a
Metropolis step whether to accept the new configuration.
It will turn out that for the detailed balance condition (4.15) we need to have
two requirements on the molecular dynamics trajectory:
\begin{itemize}
  \item Area preservation of the integration measure $D[Q]D[P]$.
  \item Reversibility of the trajectory,
        $T_{\mathrm{md}}(P',Q'|P,Q) = T_{\mathrm{md}}(-P,Q|-P',Q')$.
\end{itemize}

Reversibility can be obtained by using the so-called leapfrog integration
scheme (assuming perfect arithmetic precision): Whereas $Q$ is evolved in $n$
steps of length $\varepsilon$, the conjugate momenta start with a half-step
$\varepsilon/2$, then $(n-1)$ full steps and finally again a half-step. We
discuss the two required properties, area preservation and reversibility, for a
single such step
\begin{equation}
Q(0) \to Q(\varepsilon), \qquad
P(0) \to P(\tfrac{\varepsilon}{2}) \to P(\varepsilon),
\end{equation}
and longer trajectories are built by combining several such steps. The simplest
Euler integration scheme corresponds to the equations
\begin{equation}
\begin{split}
Q(\varepsilon) &= Q(0) + P(\tfrac{\varepsilon}{2})\,\varepsilon, \\
P(\tfrac{\varepsilon}{2}) &= P(0) - \left.\frac{\partial S}{\partial Q}\right|_{Q(0)}\frac{\varepsilon}{2}, \\
P(\varepsilon) &= P(\tfrac{\varepsilon}{2}) - \left.\frac{\partial S}{\partial Q}\right|_{Q(\varepsilon)}\frac{\varepsilon}{2}.
\end{split}
\end{equation}

Area preservation of the integration can be shown from the product of Jacobians
along this sequence\footnote{For notational convenience we use
$Q(n\varepsilon) \equiv Q_n$ and $P(n\varepsilon) \equiv P_n$.}:
\begin{equation}
\det\frac{\partial(P_1,Q_1)}{\partial(P_0,Q_0)}
= \det\frac{\partial(P_1,Q_1)}{\partial(P_{1/2},Q_1)}\,
  \det\frac{\partial(P_{1/2},Q_1)}{\partial(P_{1/2},Q_0)}\,
  \det\frac{\partial(P_{1/2},Q_0)}{\partial(P_0,Q_0)}.
\end{equation}
The right-hand side factorizes into three determinants of triangular matrices,
e.g.,
\begin{equation}
\det\frac{\partial(P_{1/2},Q_0)}{\partial(P_0,Q_0)}
= \det\begin{pmatrix} 1 & \cdots \\ 0 & 1 \end{pmatrix} = 1,
\end{equation}
and similarly for the other Jacobians. Thus the total Jacobian is 1 and the
integration measure is invariant.

In order to inspect reversibility let us first combine the Eqs.~(8.20)
starting at time 0 from some initial value $(Q_0,P_0)$. We obtain
\begin{equation}
\begin{split}
Q_1 &= Q_0 + P_0\varepsilon - \half\left.\frac{\partial S}{\partial Q}\right|_{Q_0}\varepsilon^2, \\
P_1 &= P_0 - \half\Big(\left.\frac{\partial S}{\partial Q}\right|_{Q_0}
  + \left.\frac{\partial S}{\partial Q}\right|_{Q_1}\Big)\varepsilon.
\end{split}
\end{equation}
Making a step backwards corresponds to starting at $(Q_1,P_1)$ and applying
the equations with a negative step size $-\varepsilon$. We find
\begin{equation}
\begin{split}
Q(\varepsilon - \varepsilon) = Q_0 &= Q_1 - P_1\varepsilon
  - \half\left.\frac{\partial S}{\partial Q}\right|_{Q_1}\varepsilon^2 \\
&= Q_1 - P_0\varepsilon + \half\Big(\left.\frac{\partial S}{\partial Q}\right|_{Q_0}
  + \left.\frac{\partial S}{\partial Q}\right|_{Q_1}\Big)\varepsilon^2
  - \half\left.\frac{\partial S}{\partial Q}\right|_{Q_1}\varepsilon^2 \\
&= Q_1 - P_0\varepsilon + \half\left.\frac{\partial S}{\partial Q}\right|_{Q_0}\varepsilon^2
  = Q_0, \\
P(\varepsilon - \varepsilon) = P_0 &= P_1 + \half\Big(\left.\frac{\partial S}{\partial Q}\right|_{Q_1}
  + \left.\frac{\partial S}{\partial Q}\right|_{Q_0}\Big)\varepsilon = P_0.
\end{split}
\end{equation}
Thus we end up where we have started and reversibility is established. In the
equation time always multiplies $P$ and thus this is equivalent to stating that
in the evolution $T_{\mathrm{md}}(P',Q'|P,Q) = T_{\mathrm{md}}(-P,Q|-P',Q')$, as
wanted. Actually, the sign of $P$ may be reversed, since $P$ enters quadratically
in the distributions, and thus
$T_{\mathrm{md}}(P',Q'|P,Q) = T_{\mathrm{md}}(P,Q|P',Q')$.
One may combine $n$ leapfrog steps (8.23) to build a trajectory of length
$n\varepsilon \approx 1$. This amounts to inserting a sequence of full steps
for both $Q$ and $P$ in (8.20). The discretization error for the half steps is
$\order{\varepsilon^2}$, for the normal steps $\order{\varepsilon^3}$. The
number of such steps in typical calculations is $\order{100}$.

\subsection{Completing with an accept--reject step}

Once a new configuration proposal has been obtained as the endpoint of a
leapfrog molecular dynamics trajectory, it is accepted only with a probability
depending on the change of the total Boltzmann factor. This is necessary since
the evolution of (8.18) is not exact due to the $\order{\varepsilon^2}$ errors.
The full HMC step thus consists of the following parts:
\begin{itemize}
  \item Given a configuration $Q$, generate a set of conjugate momenta $P$ from
        a Gaussian distribution $P_G(P) \propto \exp(-P^2/2)$.
  \item The molecular dynamics trajectory leads from a configuration $(P,Q)$ to
        a configuration $(P(\tau + n\varepsilon), Q(\tau + n\varepsilon))
        \equiv (P',Q')$. As discussed above, this is a deterministic process
        obeying
        \begin{equation}
        T_{\mathrm{md}}(P',Q'|P,Q) = T_{\mathrm{md}}(-P,Q|-P',Q').
        \end{equation}
  \item The new configuration is accepted with probability (see (4.16))
        \begin{equation}
        T_A(P',Q'|P,Q) = \min\!\Big(1, \frac{\exp(-H[P',Q'])}{\exp(-H[P,Q])}\Big).
        \end{equation}
\end{itemize}

We now prove that this is a valid Metropolis step in the Markov sequence of
configurations, i.e., we show detailed balance for the algorithm. The total
probability to move from $Q$ to $Q'$ results from integrating over all $P$ and
$P'$,
\begin{equation}
T(Q'|Q) = \int D[P]D[P']\,T_A(P',Q'|P,Q)\,T_{\mathrm{md}}(P',Q'|P,Q)\,e^{-P^2/2}.
\end{equation}
We then transform
\begin{equation}
\begin{split}
e^{-P^2/2-S[Q]}\,T_A(P',Q'|P,Q)
&= e^{-P^2/2-S[Q]}\min\!\Big(1,\frac{e^{-P'^2/2-S[Q']}}{e^{-P^2/2-S[Q]}}\Big)\\
&= e^{-P'^2/2-S[Q']+P^2/2+S[Q]}
   \min\!\Big(\frac{e^{-P^2/2-S[Q]}}{e^{-P'^2/2-S[Q']}},1\Big)\\
&= e^{-P'^2/2-S[Q']+P^2/2+S[Q]}\,T_A(P,Q|P',Q')\\
&= e^{-P'^2/2-S[Q']+P^2/2+S[Q]}\,T_A(-P,Q|-P',Q').
\end{split}
\end{equation}
In the last step we have used that $T_A$ is obviously even in the momentum
variables. Taking into account the reversibility (8.25) of the molecular
dynamics evolution, the integral (8.27) becomes
\begin{equation}
\begin{split}
T(Q'|Q) &= \int D[P]D[P']\,T_A(P',Q'|P,Q)\,T_{\mathrm{md}}(P',Q'|P,Q)\,e^{-P^2/2}\\
&= \int D[P]D[P']\,T_A(-P,Q|-P',Q')\,T_{\mathrm{md}}(-P,Q|-P',Q')\\
&\qquad \times e^{-S[Q']+S[Q]}\,e^{-P^2/2}\\
&= \int D[P]D[P']\,T_A(P,Q|P',Q')\,T_{\mathrm{md}}(P,Q|P',Q')\\
&\qquad \times e^{-S[Q']+S[Q]}\,e^{-P^2/2}.
\end{split}
\end{equation}
The second step is possible, since the integration measure is invariant under a
sign change of all momenta. Multiplication with $\exp(-S[Q])$ and comparing
with (8.27) then gives
\begin{equation}
e^{-S[Q]}T(Q'|Q) = e^{-S[Q']}T(Q|Q').
\end{equation}
This is the detailed balance equation (4.15) for the distribution probability
$\exp(-S[Q])$ and the proof is complete.
If the molecular dynamics equations were implemented numerically exactly, then
$H[P',Q'] = H[P,Q]$ and the acceptance probability (8.26) becomes equal to 1.
The Monte Carlo step therefore corrects for numerical errors in the
discretization. There is more to that, however. Sometimes it may be
computationally advantageous to evolve the molecular dynamics trajectory with
another, e.g., simpler action $S'$ and to accept the change using the desired
action $S$ in (8.26). This is perfectly legal. It will affect the acceptance
rate, though, which may become very small if the actions differ in any
significant manner.

\subsection{Implementing HMC for gauge fields and fermions}

Let us now apply the HMC-update to QCD with two dynamical, mass-degenerate
flavors of quarks. Using pseudofermions as in (8.8) for rewriting the
determinant, we have a system of pseudofermion fields $\varphi$ and gauge
fields $U$ distributed with the Boltzmann weight factor
\begin{equation}
\exp(-S[U]) \qquad \text{with} \qquad S[U] = S_G[U] - \varphi^\dagger(DD^\dagger)^{-1}\varphi.
\end{equation}
The pseudofermions $\varphi$ have the same indices as the Dirac field $\psi$,
i.e., color-, Dirac-, and space-time indices. The fields $\varphi$ and $U$ are
updated alternately according to the distribution (8.31).
The update of the fields $\varphi$ is simple. In fact, $\varphi$ may be easily
created by generating a complex vector $\chi$ with Gaussian distribution
$\exp(-\chi^\dagger\chi)$ and then determining $\varphi = D\chi$. The Jacobian
due to this transformation is not relevant since the gauge fields are not
changed in this step.
Once $\varphi$ has been constructed, the gauge fields $U$ are updated in a
molecular dynamics trajectory treating $\varphi$ as an external, constant
field. This leads to a new candidate gauge configuration, which is then
accepted or rejected in the final Metropolis step. This procedure of
constructing a field $\varphi$, evolving $U$, and closing with a Metropolis
step, is iterated.

We have to take into account that the gauge fields $U$ are group elements
whereas so far we have presented the molecular dynamics algorithm for variables
in flat space. Before turning to SU(3) let us discuss the simpler abelian case
of the gauge group U(1), where the link variables can be written as
$U = \exp(iQ)$. We use as conjugate variables the real number $Q$ and a
corresponding real momentum $P$. The Hamiltonian equations are identical to
(8.18). Let us consider the discretizations of the second of those,
\begin{equation}
\begin{aligned}
\dot{Q} = P \;&\to\; Q(\varepsilon) = Q(0) + P(0)\varepsilon
\;\to\; -i\ln U(\varepsilon) = -i\ln U(0) + P(0)\varepsilon\\
&\to\; U(\varepsilon) = e^{iP(0)\varepsilon}U(0).
\end{aligned}
\end{equation}
Integrating this differential equation might have led $Q$ out of the principal
value interval of the group. For U(1) this is no problem, since the relation
$U = \exp(iQ)$ automatically projects back.

Let us now implement the algorithm for the non-abelian gauge group SU(3). We
can write each link variable as, cf.\ (A.2),
\begin{equation}
U = \exp\Big(i\sum_{i=1}^{8}\omega^{(i)}T_i\Big) \equiv \exp(iQ).
\end{equation}
As for the case of U(1) our $Q$ are identified with elements of the algebra,
i.e., they are traceless hermitian matrices parametrized by the real variables
$\omega^{(i)}$. For each link variable $U_\mu(n)$ there are eight real momentum
variables $P_\mu^{(i)}(n)$ conjugate to the parameters $\omega_\mu^{(i)}(n)$,
where $n \in \Lambda$ is the space-time index. We may combine them to an
element of the algebra of SU(3) by summing them with the generators $T_i$,
\begin{equation}
P_\mu(n) = \sum_{i=1}^{8}P_\mu^{(i)}(n)\,T_i.
\end{equation}
Thus the $P_\mu(n)$ are also traceless and hermitian matrices. The Hamilton
equations involve the sum (use (A.3))
\begin{equation}
\half\sum_{n,\mu,i}\Big(P_\mu^{(i)}(n)\Big)^2 = \sum_{n,\mu}\tr\big[P_\mu(n)^2\big].
\end{equation}
We will formulate the algorithm in terms of the matrices $P_\mu(n)$ and
$U_\mu(n)$.
For the molecular dynamics evolution we also need the derivative of the action
with respect to the $Q$. Derivatives always live in the tangential space of a
manifold. For unitary groups this is the algebra (see Appendix A.1) and
derivatives thus are along some path in the algebra. The derivative of some
function $f(U)$ of a group element in the algebra direction $T_i$ may be
defined as,\footnote{This definition can be understood from differentiating
$\ln(U) = i\sum_k \omega^{(k)}T_k$ with respect to $\omega^{(i)}$. We remark
that this definition is not unique, but the outcome of (8.36) is, if $f(U)$ is
a class function.}
\begin{equation}
\nabla^{(i)}f(U) \equiv \frac{\partial f(U)}{\partial\omega^{(i)}}
= \left.\frac{\partial}{\partial\omega^{(i)}}f\big(e^{i\omega^{(i)}T_i}U\big)\right|_{\omega=0}.
\end{equation}
In this way we find, e.g., ($A, B \in \mathrm{SU}(3)$)
\begin{equation}
\nabla^{(i)}(AUB) = iAT_iUB, \qquad
\nabla^{(i)}(AUUB) = iAT_iUUB + iAUT_iUB,
\end{equation}
and so on. The leapfrog discretized molecular dynamics equations (8.20) for
altogether $n$ steps then become:
\begin{enumerate}
  \item \textbf{Pseudofermions}\\
        Generate the pseudofermion field $\varphi = D\chi$, where $\chi$ is
        distributed according to $\exp(-\chi^\dagger\chi)$.
  \item \textbf{Conjugate fields}\\
        Given a gauge configuration $U_0$ generate $P_0$ (eight real numbers
        for each $(n,\mu)$) according to the Gaussian distribution
        $\exp(-\tr[P^2])$.
  \item \textbf{Initial step}\\
        $P_{\frac12} = P_0 - \frac{\varepsilon}{2}F[U,\varphi]\big|_{U_0}$.
  \item \textbf{Intermediate steps}\\
        Full steps for $k = 1, \ldots, n-1$:
        $U_k = \exp(i\varepsilon P_{k-\frac12})U_{k-1}$, \qquad
        $P_{k+\frac12} = P_{k-\frac12} - \varepsilon F[U,\varphi]\big|_{U_k}$.
  \item \textbf{Final step}\\
        $U_n = \exp(i\varepsilon P_{n-\frac12})U_{n-1}$, \qquad
        $P_n = P_{n-\frac12} - \frac{\varepsilon}{2}F[U,\varphi]\big|_{U_n}$.
  \item \textbf{Monte Carlo step}\\
        Accept--reject step at the end of the trajectory: accept the changes if
        a random number $r \in [0,1)$ is smaller than
        \begin{equation}
        \begin{aligned}
        &\min\Big(1,\;\exp\big[\tr P^2 - \tr P'^2 + S_G[U] - S_G[U']\\
        &\qquad + \varphi^\dagger\big((DD^\dagger)^{-1} - (D'D'^\dagger)^{-1}\big)\varphi\big]\Big).
        \end{aligned}
        \end{equation}
\end{enumerate}
In these equations
\begin{equation}
F[U,\varphi] = \sum_{i=1}^{8}T_i\nabla^{(i)}\Big(S_G[U] + \varphi^\dagger(DD^\dagger)^{-1}\varphi\Big) \in su(3)
\end{equation}
denotes the driving force, an element of the algebra $su(3)$.
For the pseudofermion contribution in the accept--reject step one re-uses the
vector that has been computed for the force in step (5), i.e.,
$(D'D'^\dagger)^{-1}\varphi$.
A technical problem is the exponentiation $\exp(i\varepsilon P)$ in step (4)
which maps an element of the algebra to a group element. This may be done
either by a series expansion or, faster, by utilizing the Cayley--Hamilton
relation as done in [9].

The force term (8.39) involves derivatives with regard to the algebra elements.
As an example let us determine the force due to the contribution from the
Wilson gauge action (4.20) involving a particular link $U$. It is given by
\begin{equation}
-\frac{\beta}{6}\tr\big[UA + A^\dagger U^\dagger\big],
\end{equation}
where $A$ denotes the sum of adjacent staples (4.20). The derivative of the
trace-term then is
\begin{equation}
\nabla^{(i)}\tr\big[UA + A^\dagger U^\dagger\big]
= \tr\big[iT_iUA - iA^\dagger U^\dagger T_i\big]
= \tr\big[iT_i\big(UA - A^\dagger U^\dagger\big)\big].
\end{equation}
The resulting force is
\begin{equation}
-\frac{\beta}{6}\sum_{i=1}^{8}T_i\,\tr\big[iT_i\big(UA - A^\dagger U^\dagger\big)\big]
= -\frac{\beta}{12}\,i\big(UA - A^\dagger U^\dagger\big).
\end{equation}
We have used that $i(UA - A^\dagger U^\dagger)$ is traceless and hermitian and
thus is a linear combination of the $T_j$. For such a sum one has the identity
\begin{equation}
\sum_j T_j\,\tr\Big[T_j\sum_k c_k T_k\Big] = \half\sum_j c_j T_j.
\end{equation}
The trace projects out the contribution of $T_j$ (introducing a factor of
1/2), and the sum reconstructs the expression.
The contribution of the fermion action to the force requires the evaluation of
(here we use $\partial M^{-1}/\partial\omega \allowbreak = \allowbreak -M^{-1}(\partial M/\partial\omega)M^{-1}$)
\begin{equation}
\begin{split}
\nabla^{(i)}\varphi^\dagger(DD^\dagger)^{-1}\varphi
&= -\varphi^\dagger(DD^\dagger)^{-1}\big[\nabla^{(i)}(DD^\dagger)\big](DD^\dagger)^{-1}\varphi \\
&= -\Big((DD^\dagger)^{-1}\varphi\Big)^\dagger
   \Big[D\frac{\partial D^\dagger}{\partial\omega^{(i)}}
   + \frac{\partial D}{\partial\omega^{(i)}}D^\dagger\Big](DD^\dagger)^{-1}\varphi .
\end{split}
\end{equation}
In each substep of the trajectory one therefore has to compute the derivative
of $D$ and the inverse of $DD^\dagger$, which is the time-consuming part of the
algorithm.
In computing the derivative of the Dirac operator one has to observe
non-commutativity of the group elements. Some Dirac operators involve
contributions, where a link variable occurs more than once. An obvious example
is the overlap operator in a rational or series approximation (see
Sect.~7.4.3). In such a case one has to be particularly careful in determining
the derivative.
The Wilson Dirac operator (5.51) is simply linear in the gauge field variables
and in this case computing the derivative is straightforward. We find for the
contribution of a particular link $U_\mu(k)$
\begin{equation}
\frac{\partial D(n|m)}{\partial\omega_\mu^{(i)}(k)}
= -i\frac{1+\gamma_\mu}{2a}T_iU_\mu(k)\delta_{n+\hat{\mu},m}\delta_{n,k}
  + i\frac{1-\gamma_\mu}{2a}U_\mu(k)^\dagger T_i\delta_{n-\hat{\mu},m}\delta_{m,k}.
\end{equation}
Due to the Metropolis decision at the end of each trajectory the HMC algorithm
is exact. Any numerical error done in the discretization is ``repaired'' by the
probabilistic accept--reject step. However, the acceptance rate may become low.
One therefore has to adjust the step size $\varepsilon$ to prevent too large
aberrations. Typical numbers are 50--100 steps for a complete trajectory with
an integrated time of $n\varepsilon \approx 1$. The estimated computational
effort for the HMC algorithms grows only like $\propto V^{5/4}$ [10, 11].

\section{Other algorithmic ideas}

\subsection{The R-algorithm}

The HMC algorithm is exact. However, in the way presented here it is applicable
only for positive definite fermion actions (like the overlap action) or for
duplicated number of fermion species.\footnote{Suggestions how to include
single fermion species will be discussed in Sect.~8.3.3.} It evolved out of
earlier updating methods, which also used microcanonical (or Langevin) type
differential equations, without the correcting Monte Carlo step.
In all of the so-called hybrid algorithms [7] the system follows a molecular
dynamics trajectory with refreshing the conjugate momenta regularly. This
process is equivalent to an evolution according to a stochastic differential
equation. It has been argued that systems out of equilibrium will be driven
toward equilibrium and stay there during the updating process. The crucial
problem is the necessary discretization, which introduces systematic errors in
some order of the discretization time step.
One such method, the so-called R-algorithm [12], is advocated for
Kogut--Susskind or staggered fermions (cf.\ Chap.~10). The central idea is to
use the effective action discussed in (8.10). The molecular dynamics updating
follows steps like those discussed earlier. The change of the effective action
is estimated stochastically [13] with the help of noisy pseudofermions. The
change of the effective fermion action due to a change of the gauge variable
then is
\begin{equation}
\begin{split}
\Delta S_F^{\mathrm{eff}} &= S_F^{\mathrm{eff}}[U_\mu(n)'] - S_F^{\mathrm{eff}}[U_\mu(n)]\\
&= \sum_{n,\mu,i}\frac{\partial S_F^{\mathrm{eff}}[U_\mu(n)]}{\partial\omega_\mu^{(i)}(n)}
   \Delta\omega_\mu^{(i)}(n)
   + \order{\Big(\Delta\omega_\mu^{(i)}(n)\Big)^2}.
\end{split}
\end{equation}
For the derivative of the effective action we obtain
\begin{equation}
\frac{\partial S_F^{\mathrm{eff}}}{\partial\omega_\mu^{(i)}(n)}
= -\frac{\partial\tr[\ln(D[U])]}{\partial\omega_\mu^{(i)}(n)}
= -\tr\Big[D[U]^{-1}\frac{\partial D[U]}{\partial\omega_\mu^{(i)}(n)}\Big].
\end{equation}
Note that the arguments are matrices and these are usually not commutative.
However, the cyclicity property of the trace simplifies the expression. For
actions where the link variables occur linearly, the derivative of $D$ is
particularly simple.

As there is no correction MC step, the method will show a dependence on the
chosen value of the discretization time step $\varepsilon$. By using clever
modifications the errors can be made $\order{\varepsilon^2}$. This dependence
has to be dealt with by repeating the computations for different values of
$\varepsilon$ and extrapolating to 0.
At each time step one computes pseudofermions $\varphi = M^\dagger\chi$ from
Gaussian distributed random vectors $\chi$. More details can be found in [12].

\subsection{Partial updates}

For updating methods that are based on evolution equations (like the HMC and
the R-algorithm or any such method working with Langevin-type differential
equations) one has to calculate the force term. This becomes quite involved for
Ginsparg--Wilson-type Dirac operators. There one has to use computationally
costly algorithms already for evaluating $D$ itself.
It is therefore tempting to shortcut the update by updating larger chunks of
gauge variables at once. However, as has been argued earlier, this may run into
acceptance problems when working at large volumes: A naive implementation will
lead to a numerical cost of $\order{V^2}$. One therefore has to invest some
effort in providing better suggestions for the change of gauge variables,
similar to the ideas originally leading to the hybrid algorithms.
It has been advocated [14, 15] to interleave gauge field updates on part of the
lattice based on a local gauge action that approximates in some way the effect
of the fermionic action. This may be done, e.g., by considering a set of closed
gauge loops with coefficients properly tuned. This partial-global change of
variables then undergoes an accept/reject step, where the ratio of determinants
is estimated stochastically. Since the change of the effective action may be
large, one has to be particularly careful in getting a reliable estimate of
that ratio. One therefore invokes techniques like polynomial approximations,
preconditioning, and multi-pseudofermions.
For real, positive determinants it is then possible to simulate single fermion
species. This is true for Ginsparg--Wilson operators, which are normal matrices
with eigenvalues that have positive real part. For such operators one may use a
pseudofermion action of the form $(\varphi, D^{-1}\varphi)$ and Gaussian
$\chi$ with $\varphi = \sqrt{D}\,\chi$ to generate the correct distribution. For
taking the square root one uses methods like that discussed for the
construction of the overlap operator.

\subsection{Polynomial and rational HMC}

In the standard HMC algorithm the Dirac operator enters quadratically as
$DD^\dagger$ to guarantee positivity of the real part of all eigenvalues
necessary for the convergence of the pseudofermion Gaussian integral. Thus in
that form only even numbers of fermions can be simulated.
If $\det[D]$ is positive, there is a way to circumvent this restriction. One
approximates $D^{-1}$ by an operator $TT^\dagger$ for the HMC trajectory
calculations and corrects for the approximation in the accept--reject step. The
approximations suggested are either by polynomials [16--18] or by rational
functions [19--22], in the spirit of the evaluations of the overlap operator in
Chap.~7.
As an example we discuss the suggestions of [18]. The inverse of the Dirac
operator is approximated by an even polynomial,
\begin{equation}
D^{-1} \approx P_{2n}(D) = \prod_{k=1}^{n}(D - z_{2k-1})(D - z_{2k-1}^*),
\end{equation}
with $z_k = 1 - \exp(i\,2\pi k/(2n+1))$. Assuming $\gamma_5$-hermiticity one
has
\begin{equation}
\det[D - z_k^*] = \det[\gamma_5 D^\dagger\gamma_5 - z_k^*] = \det[D - z_k]^*,
\end{equation}
and we therefore may split $P_{2n}$ into a product
\begin{equation}
P_{2n}(D) = T_n(D)T_n^\dagger(D) \qquad \text{with} \qquad
T_n(D) = \prod_{k=1}^{n}(D - z_{2k-1}).
\end{equation}
The correct determinant is
\begin{equation}
\det[D] = C\det[T_n^\dagger T_n]^{-1} \qquad \text{with} \qquad
C = \det[DT_n^\dagger T_n].
\end{equation}
The HMC trajectory is then computed using a pseudofermion action
\begin{equation}
S = -\varphi^\dagger T_n^\dagger(D)T_n(D)\varphi,
\end{equation}
thus providing for the second factor in (8.51), whereas the correction factor
$C$ is dealt with in the Metropolis accept--reject step at the end of the
trajectory. Such methods allow one to work with odd numbers of quark flavors
using the exact MC algorithm.

\subsection{Multi-pseudofermions and UV-filtering}

In the accept--reject Metropolis step, as it is necessary at the end of each
HMC trajectory or in (partial) global algorithms, the stochastic determination
with noisy pseudofermion estimation introduces additional noise. Although
formally correct, the induced fluctuations of the fermionic force may
destabilize the algorithm [2, 20, 23], enforcing smaller step sizes.
A possible way to reduce the noise of the stochastic estimator is to use more
than one set of pseudofermion variables. Assuming positivity for $M$ one uses
the identity
\begin{equation}
\det[M] = \Big(\det[M^{1/n}]\Big)^n.
\end{equation}
The stochastic estimator then is
\begin{equation}
\det[M] \propto \prod_{k=1}^{n}\int D[\varphi_R]D[\varphi_I]\,e^{-\varphi^\dagger M^{-1/n}\varphi},
\end{equation}
where $M$ is either a positive definite Dirac operator or $M = DD^\dagger$. The
matrix $M^{-1/n}$ can be approximated by a rational function or a polynomial in
$M$.
A variant of this method [24] splits the determinant
\begin{equation}
\det[M] = \det[M']\,\det[M'^{-1}M]
\end{equation}
with separate pseudofermion fields for both factors. By choosing different
parameters for the matrices $M$ and $M'$ one may improve the updating
performance; the step size can be increased while the acceptance rate remains
stable. In the applications one chooses different mass parameters for the two
components; thus the method has also been called mass preconditioning.
Another method to improve the estimate of the ratio of determinants is to
introduce an UV-filter [25]. The idea is to reduce the spread of the
eigenvalues of the operator to be stochastically estimated. One defines a
reduced matrix
\begin{equation}
D_{m,r} = D\,e^{f(D)},
\end{equation}
where $f(D)$ is chosen to be a suitable low-order polynomial in $D$ with real
coefficients such that the eigenvalues of $D_{m,r}$ are concentrated around
$z = 1$ in the complex plane, i.e., $e^{f(D)}$ roughly approximates $D^{-1}$. In
the acceptance step one then computes (using $\gamma_5$-hermiticity)
\begin{equation}
\frac{\det[D^\dagger D]_{\mathrm{new}}}{\det[D^\dagger D]_{\mathrm{old}}}
= \exp(-2\tr[f(D)_{\mathrm{new}}] + 2\tr[f(D)_{\mathrm{old}}])\,
  \frac{\det[D_{m,r}^\dagger D_{m,r}]_{\mathrm{new}}}{\det[D_{m,r}^\dagger D_{m,r}]_{\mathrm{old}}}.
\end{equation}
The computation of the trace over polynomials of $D$ is possible for simple
Dirac operators but may be an obstacle for extended ones.

\subsection{Further developments}

Both the multi-pseudofermion method and the polynomial approximation are in
spirit related to the earlier suggestion of a multiboson algorithm [26].
Depending on the Dirac operator there are also various methods of
preconditioning (see also Sect.~6.2.4). Mass preconditioning and variants
thereof are one possible improvement. For operators involving only nearest
neighbor couplings one may split the lattice into even and odd sites and in
this way simplify the calculations. This method is called even-odd
preconditioning.
Another hierarchically organized approach splits the lattice into nonoverlapping
domains and solves the system on the sub-domains. This so-called Schwarz
alternating procedure [27--29] has been tested for both preconditioning and
dynamical fermion algorithms. It is a domain decomposition method for solving
elliptic differential equations (introduced already in the nineteenth century).
Although at least for HMC algorithms one expects the necessary computer time to
grow with $V^{5/4}$ for a given set of couplings, the growth as a function of
the lattice spacing or the quark mass is not as clear. Since we expect the
quark mass to be proportional to the mass squared of the lightest pseudoscalar,
it is useful to discuss the computational cost in such units. The expected
functional form for the numerical cost is [30, 31]
\begin{equation}
\mathrm{cost} \propto L^{z_L}\,a^{-z_a}\,M_\pi^{-z_\pi}.
\end{equation}
Various aspects enter such a consideration. As mentioned the volume scales with
$z_L \approx 5$. The condition number of the Dirac operator matrix decreases
with the pion mass, thus the CGR solver needs more iterations. The step size in
the HMC trajectory will have to be decreased and the autocorrelation time may
increase. The various studies with, e.g., the Wilson fermion action have led to
values of [30] $z_a \approx 7$ and $z_\pi \approx 6$, but also smaller values
$\order{3-4}$ have been quoted [31].
The field of algorithms for dynamical fermions has been evolving over the years
and the development has not yet ceased. Various ideas have been presented,
which are not discussed here. Other ideas are certain to come and we advise to
check the proceedings of the yearly lattice conferences.

\section{Other techniques using pseudofermions}

In Sect.~8.1.3 we have remarked that pseudofermion fields are used in noisy
estimator techniques designed for evaluating fermionic observables. The
starting point is the structural similarity between bosonic and fermionic
generating functionals, (5.32) and (8.7).
Let us start discussing the idea of noisy estimators by using an example of
high importance: the value of the fermion condensate $\ev{\bar{\psi}\psi}$. In
Chap.~7 we have found that this is essentially given by $\tr[D^{-1}]$.
Straightforward calculation of the trace requires the determination of
propagators for $12|\Lambda|$ point sources, clearly an impossible task. For
the noisy estimator approach we consider a set of uncorrelated complex random
numbers $\chi_i \equiv \chi_{R,i} + i\chi_{I,i}$ with unit variance,
\begin{equation}
\ev{\chi_i\chi_j^*}_\chi = \delta_{ij}.
\end{equation}
For example, we may choose for the $\chi_i$ the Gaussian distribution
\begin{equation}
\ev{\cdots}_\chi = \pi^{-N}\int D[\chi_R]D[\chi_I]\exp(-\chi^\dagger\chi)\cdots.
\end{equation}
We may approximate this integral by an average over $K$ random vectors
$\chi^{(k)}$ with entries $\chi_i^{(k)}$ distributed according to the
probability density (8.60),
\begin{equation}
\delta_{ij} = \ev{\chi_i\chi_j^*}_\chi \approx \frac{1}{K}\sum_{k=1}^{K}\chi_i^{(k)}\chi_j^{(k)*}.
\end{equation}
Using this one can estimate the trace
\begin{equation}
\tr[D^{-1}] = \sum_{i,j}(D^{-1})_{ji}\delta_{ij}
\approx \frac{1}{K}\sum_{k=1}^{K}\chi^{(k)\dagger}D^{-1}\chi^{(k)}.
\end{equation}
Again the random variables agree in number with the fermionic degrees of
freedom. The number $K$ of noise vectors considered depends on the desired
accuracy of the estimate.
Gaussian noisy pseudofermions may also be utilized to evaluate the change of
the fermionic effective action [32], which is required for the Monte Carlo
update algorithm. For matrices $M$ with strictly positive real part of the
eigenvalues we have
\begin{equation}
\begin{split}
\det[M] &= \frac{\int D[\chi_R]D[\chi_I]\exp(-\chi^\dagger\chi)\exp\big(-\chi^\dagger(M^{-1}-1)\chi\big)}
        {\int D[\chi_R]D[\chi_I]\exp(-\chi^\dagger\chi)}\\
&= \Big\langle\exp\big(-\chi^\dagger(M^{-1}-1)\chi\big)\Big\rangle_\chi .
\end{split}
\end{equation}
This identity follows from (8.7) and may be used to determine directly the
change of the effective 2-flavor action:
\begin{equation}
\begin{split}
\exp(-\delta S_{F,\mathrm{eff}}) &= \det[(D'D'^\dagger)(DD^\dagger)^{-1}]
  = \det[D'^\dagger(DD^\dagger)^{-1}D']\\
&= \Big\langle\exp\big(-\chi^\dagger\big(D'^{-1}DD^\dagger(D'^\dagger)^{-1}-1\big)\chi\big)\Big\rangle_\chi\\
&= \Big\langle\exp\big(-\chi^\dagger D'^{-1}\big(DD^\dagger - D'D'^\dagger\big)(D'^\dagger)^{-1}\chi\big)\Big\rangle_\chi .
\end{split}
\end{equation}
Again $\ev{\cdots}_\chi$ is approximated by a finite number of noise vectors.
It has been suggested in [33] to replace the sum over many $\varphi$-ensembles
by just one individual for the Monte Carlo acceptance decision. Determining
that ratio for the updating step either way is correct in the average. However,
this type of approach involves a Monte Carlo sum within a Monte Carlo sum and,
as has been discussed in, e.g., [34], corresponds to an updating algorithm where
the computer time grows with the square of the lattice volume at least, not
considering other effects like critical slowing down.
Whenever one wants a precise calculation of quark propagators, one uses
numerical inverters. However, it is expensive to compute the inverse matrix for
many different sources. This is, e.g., necessary for sums over disconnected
pieces (see Fig.~6.1). In such a case one might be willing to trade lower
accuracy for a larger number of source points. Noisy pseudofermions provide a
tool for approximating propagator elements. Exploiting (8.7) one may compute
them as expectation values of $\varphi_n\varphi_k^*$ in a bosonic Gaussian
distribution:
\begin{equation}
D_{nm}^{-1} = \Big[(D^\dagger D)^{-1}\Big]_{nk}D_{km}^\dagger
= \frac{\int D[\varphi_R]D[\varphi_I]\,e^{-\varphi^\dagger D^\dagger D\varphi}\,\varphi_n\varphi_k^*}
       {\int D[\varphi_R]D[\varphi_I]\,e^{-\varphi^\dagger D^\dagger D\varphi}}\,D_{km}^\dagger .
\end{equation}
As long as $D$ has nonvanishing eigenvalues this Gaussian integral is
well-defined. It can be evaluated using a stochastic process for the
$\varphi$.
Formulas of this type have also been used in early suggestions to implement
dynamical fermions [35] for computing the expressions necessary for the change
of the effective fermionic action in (8.46).

\section{The coupling-mass phase diagram}

\subsection{Continuum limit and phase transitions}

In Sect.~3.5.4 we discussed the true continuum limit for pure gauge theory. We
used the renormalization group analysis result (3.88) for the lattice spacing
$a$ as a function of the inverse coupling $\beta$ to identify the limit
$\beta \to \infty$ as the limit where $a \to 0$ and the continuum physics is
recovered. Here we address the continuum limit a second time, but now
reformulate it using statistical mechanics language.
To sketch the idea, we first consider the simple familiar 2D Ising model.
Without external magnetic field, there is only a single coupling parameter.
Depending on its value one finds two phases, a disordered ``hot'' phase and a
phase with a nonvanishing magnetization, the ``ordered'' phase. A second-order
phase transition separates the two phases. The phase transition occurs at a
specific value of the coupling, the so-called critical point where the
correlation length $\xi$ of the system diverges. The correlation length can be
defined through the connected correlation function of two spin variables
located at the origin and some site $n$. For large distances between the spins
the correlation function decays exponentially
\begin{equation}
\ev{s_ns_0} - \ev{s_n}\ev{s_0} \propto \exp\Big(-\frac{|n|}{\xi}\Big)
\qquad \text{for } n \to \infty.
\end{equation}
The correlation length $\xi$ is a measure for the distance over which two spins
are correlated. When the coupling is driven close to its critical value $\xi$
becomes large and the spins are correlated over long distances. The local
structures extending only over sites close to each other become less and less
important.
Why is the notion of a diverging correlation length relevant for lattice field
theory? We have learned that the physical mass $M_H$ of some hadron can be
computed from the two-point function of a hadron interpolator $O_H$,
\begin{equation}
\ev{O_H(n_t)\,O_H(0)^\dagger} \propto e^{-n_t a M_H} = e^{-n_t/\xi_H}.
\end{equation}
Here we have defined the correlation length $\xi_H$ for the interpolator $O_H$
as
\begin{equation}
\xi_H = \frac{1}{aM_H}.
\end{equation}
We stress that interpolators for different hadrons will in general give rise to
different correlation lengths proportional to each other. Since $M_H$ is
constant all correlation lengths $\xi_H$ of physically observable states should
diverge in the continuum limit, $a \to 0$. As in the systems of statistical
mechanics we must drive our couplings, the inverse gauge coupling $\beta$ and
the quark mass $m$, toward critical values where the correlation length
diverges. At these critical values we can construct a continuum limit. It is
therefore of utmost importance to identify such critical points in the phase
diagram.

\subsection{The phase diagram for Wilson fermions}

In QCD we expect that the correct continuum limit can be obtained at vanishing
gauge coupling $g = 0$, corresponding to $\beta = \infty$, and a vanishing
fermion mass parameter. However, also for other values of these couplings we
have transitions separating phases of different physics.
For Wilson fermions it is usual to discuss the phase diagram in terms of the
inverse gauge coupling $\beta = 6/g^2$ and the hopping parameter $\kappa$
introduced in (5.55). The hopping parameter is related to the bare quark mass
in the continuum limit through $\kappa = 1/(2am + 8)$. In the hopping parameter
notation the Dirac operator is given by (5.57).
When discussing a phase diagram, it is always useful to first consider simple
limiting cases. For $\kappa = 0$ the hopping term is turned off and the Dirac
operator becomes a unit matrix. This limit corresponds to infinitely heavy
quarks. This is the situation of the pure gauge theory (the quenched case)
discussed already earlier. In that limit the string tension defined through the
Wilson loop may be interpreted as order parameter. As we have discussed in
Sect.~3.4.1, at $\beta = 0$ (infinite bare gauge coupling $g$) the theory is
confining and Monte Carlo evidence is that this confinement phase extends over
the whole range of values toward $\beta \to \infty$ ($g = 0$). Actually, since
Monte Carlo calculations are limited, all we can really say is that for the
pure gauge theory there is no numerical evidence for a phase transition to a
deconfined phase for large enough lattices.
Another limiting case is that of free lattice fermions. This is obtained for
$\beta = \infty$ ($g = 0$), where all plaquettes approach unity. On an infinite
lattice this is equivalent to setting all link variables to 1. With the help of
lattice Fourier transformation one can compute the quark propagators for the
free case and find that the correlation length is given by the inverse of the
quark mass. Thus the correlation length diverges for vanishing quark mass
$m = 0$, which corresponds to a value of $\kappa = 1/8$.
For finite values of $\beta$ one finds that the smallest observed mass in the
system (the mass of the pion) vanishes at some value $\kappa$ which we denote
by $\kappa_{\mathrm{ch}}(\beta)$. Since we associate the pion with the
Goldstone mode of spontaneous chiral symmetry breaking, this is the point where
at finite $\beta$ the quark mass vanishes. In this indirect way we get our
first information on the quark mass itself. This line of
$\kappa = \kappa_{\mathrm{ch}}(\beta)$ appears to continue from
$\kappa_{\mathrm{ch}}(\beta = \infty) = 1/8$ down to
$\kappa_{\mathrm{ch}}(\beta = 0)$. There are strong coupling expansion results
for meson masses [36, 37] indicating that
$\kappa_{\mathrm{ch}}(\beta = 0) \approx 1/4$. The resulting phase diagram is
sketched in Fig.~8.1.

\begin{figure}[htbp]
  \centering
  \includegraphics[width=0.8\textwidth]{fig81.png}
  \caption*{Fig. 8.1. \emph{The phase diagram for QCD with Wilson fermions. The
  coupling constants are $\beta = 6/g^2$ and the hopping parameter $\kappa$. The
  dashed curves are lines of constant mass ratios of, e.g.,
  $M_\pi/M_\rho = \mathrm{const}$.}}
  \label{fig:81}
\end{figure}

At a given $\beta$ the quark masses thus assume values between infinity and 0
when $\kappa$ runs from 0 to $\kappa_{\mathrm{ch}}(\beta)$. Driving $\kappa$
towards $\kappa_{\mathrm{ch}}(\beta)$ the pion mass decreases toward 0, but
other masses of the system, like the $\rho$ or the nucleon mass, approach
nonzero values. At some value of $\kappa$ the mass ratio of, e.g.,
$M_\pi/M_\rho$ will assume its experimental value. At this value of $\kappa$ we
then use the dimensionless combinations $aM$ and the experimental values $M$ to
determine the lattice spacing $a$ in order to set the scale.
The collection of such points for varying $\beta$ then defines a ``physical''
curve (dashed curves in Fig.~8.1). Along such a curve the lattice spacing
decreases toward zero for increasing $\beta$ and the system approaches the
continuum limit. We emphasize that the dynamics of the sea quarks of course
also influences the lattice spacing $a$, which thus is a function of both
coupling parameters $a = a(\beta, \kappa)$. Curves of constant lattice spacing
thus will not simply be lines of constant $\beta$. Also we point out that the
lattice spacing is nonzero along the chiral curve
$\kappa = \kappa_{\mathrm{ch}}(\beta < \infty)$.\footnote{The situation is
similar to that of spontaneous symmetry breaking in the $N$-component
$\varphi^4$-theory. There the two couplings are the nearest neighbor-coupling
$g$ (playing the role of our gauge coupling) and an external magnetic field $h$
(related to the fermion mass in QCD). In the broken phase, for $g > g_{\mathrm{crit}}$
and $h = 0$ the order parameter $\ev{\varphi}$ is nonzero and there are
Goldstone modes as well as massive modes. The lattice spacing is zero only
toward the critical point $g \to g_{\mathrm{crit}}$. When the external field is
switched on the Goldstone modes become massive, similar to the pions.}
Approaching the continuum limit along the chiral curve is another alternative
suitable for mass-independent renormalization schemes. In this situation the
lattice spacing $a$ for given gauge coupling is defined in the chiral limit,
i.e., on the chiral curve.
We have to add an important remark to the above discussion. As will be detailed
in Chap.~12, by studying lattices with finite time extent (in physical units)
one may simulate QCD at finite temperature. This introduces another parameter,
the temperature, giving rise to a more complicated phase diagram. One finds a
phase transition to a deconfined phase at some critical temperature, i.e., if
the time (or space) extent becomes sufficiently small. Here we only stress that
the above discussion of the phase structure is based on the assumption that the
time extent of our lattice is large enough such that we always remain in the
confined phase.

\subsection{Ginsparg--Wilson fermions}

The range of the bare mass parameter $am$ for a Ginsparg--Wilson Dirac operator
in the form of (7.41) is limited to $[0, 2]$. The value $am = 0$ defines the
point of massless fermions and (spontaneously broken) chiral symmetry. At the
maximum value $am = 2$ the fermions decouple from the system; this corresponds
to the limit of infinitely heavy quarks.
Since the curve of chiral symmetry is just a line at $m = 0$ and does not depend
on the gauge coupling (unlike the discussed case of the Wilson action), the
phase diagram looks simpler than that of Fig.~8.1. The variable $am$ replaces
$\kappa$ and ranges from $am = 2$ for the pure gauge theory (quenched limit) to
$m = 0$ for the chiral limit. Again one can find curves of constant mass ratios
approaching the continuum limit at $\beta \to \infty$, $m = 0$.
Let us address another aspect of using a Dirac operator which for vanishing
quark mass obeys the Ginsparg--Wilson equation. For such a Dirac operator the
eigenvalue spectrum is restricted to a circle in the complex plane. The
smallest distance of this circle to the origin is given by the bare quark mass
and the smallest possible eigenvalue is given by this mass. For nonvanishing
mass the inverse Dirac operator, i.e., the quark propagator can always be
computed, and (in principle) the chiral line can be approached.
For non-chiral Dirac operators, e.g., the Wilson Dirac operator, one has no such
restriction of the eigenvalues and they fluctuate depending on the gauge field.
Thus even for nonvanishing quark mass zero eigenvalues of the Dirac operator may
show up for exceptional configurations (compare Sect.~6.2.5). Consequently the
quark propagator cannot be computed for such configurations. Thus the
fluctuations of the Dirac eigenvalues often forbid a close approach to the
chiral line $\kappa_{\mathrm{ch}}(\beta)$. The situation improves by using
larger values of $\beta$, larger lattices, or more advanced Dirac operators
with smaller fluctuations of their eigenvalues.

\section{Full QCD calculations}

In full QCD calculations the large mass quarks $c$, $b$, and $t$ contribute
negligibly to the dynamics of the quark sea. It is therefore an excellent
approximation to consider only two mass-degenerate light quarks $u$, $d$ and
one heavier quark $s$ for the sea. Even that poses an algorithmic problem. For
the non-exact R-type algorithms discussed in Sect.~8.3.1 the number of quark
species of a kind may be easily fixed by a pre-factor; however the flavor
mixing and the continuum limit are disputable. For the exact HMC-algorithm in
its original form, on the other hand, positivity requires an even number of
mass-degenerate fermions. There are, however, variants where odd numbers of
fermions may be implemented with the help of polynomial approximations as
mentioned in Sect.~8.3.
There are several international collaborations studying dynamical fermions for
various types of Dirac operators. If we tried to mention them explicitly this
would imply omitting some of them even at the time of publication of this text.
We therefore just remark that for the simple improved Wilson action spatial
lattice sizes up to 50--100 are presently achieved for two light flavors and
some even including the strange quark. Similar lattice sizes are being studied
for twisted mass and staggered fermions (see Chap.~10). For the numerically
more demanding domain wall formulation (Chap.~10) one typically works on
lattices half that size (in lattice units). This roughly reflects the increased
numerical cost, i.e., one to two orders of magnitude. First results for overlap
fermions on smaller lattices, again a factor of two smaller than the domain
wall sizes, are being produced, although with certain restrictions like fixed
topological sectors due to algorithmic problems for the tunneling between
different sectors. In most of these collaborations a wide spectrum of physics
questions is attacked by the participating scientists.

\begin{figure}[htbp]
  \centering
  \includegraphics[width=0.9\textwidth]{fig82.png}
  \caption*{Fig. 8.2. \emph{Spectroscopy results from a fully dynamical
  simulation by the Budapest--Marseille--Wuppertal collaboration [38]. (From
  S.~D\"urr et al.: Science 322, 1224 (2008). Reprinted with permission from
  AAAS)}}
  \label{fig:82}
\end{figure}

To give one example of a calculation with dynamical fermions we show in
Fig.~8.2 the result of a spectroscopy calculation with two light dynamical
fermions and a dynamical strange quark from the Budapest--Marseille--Wuppertal
collaboration [38]. The simulation was done with the improved (compare
Chap.~9) Wilson gauge- and Wilson fermion action on lattices with a spatial
extent up to 4 fm. The gauge fields entering the Dirac operator were treated
with six levels of stout smearing. Pion masses down to 190 MeV were reached and
finite volume effects, as well as effects of resonances, were taken into
account in the analysis. In Fig.~8.2 the masses for different mesons and
baryons are shown, extrapolated to physical quark masses and vanishing lattice
spacing.
Comparing these results with the quenched spectroscopy calculation of Fig.~6.5
one finds that the discrepancies between the lattice results and the
experimental numbers have vanished. Figure 8.2 demonstrates that when taking
into account the effects of light dynamical quarks, for a wide range of
different quantum numbers a controlled \emph{ab initio} lattice spectroscopy
calculation leads to excellent agreement between numerical results and the
experimental data.

\section*{References}

\begin{enumerate}[label={[\arabic*]}]
  \item D.~H.~Weingarten and D.~N.~Petcher: Phys.\ Lett.\ B 99, 333 (1981) 188
  \item A.~Alexandru and A.~Hasenfratz: Phys.\ Rev.\ D 66, 094502 (2002) 190, 201
  \item D.~J.~E.~Callaway and A.~Rahman: Phys.\ Rev.\ D 28, 1506 (1983) 192
  \item J.~Polonyi and H.~W.~Wyld: Phys.\ Rev.\ Lett.\ 51, 2257 (1983) 192
  \item G.~Parisi and Y.-S.~Wu: Sci.\ Sin.\ 24, 483 (1981) 192
  \item J.~R.~Klauder: in \emph{Recent Developments in High-Energy Physics};
        Acta Phys.\ Austriaca Suppl.\ 25, edited by H.~Mitter and C.~B.~Lang,
        p.~251 (Springer, Wien, New York 1983) 192
  \item S.~Duane and J.~B.~Kogut: Phys.\ Rev.\ Lett.\ 55, 2774 (1985) 192, 199
  \item S.~Duane and J.~B.~Kogut: Nucl.\ Phys.\ B 275, 398 (1986) 192
  \item C.~Morningstar and M.~Peardon: Phys.\ Rev.\ D 69, 054501 (2004) 198
  \item R.~Gupta, G.~W.~Kilcup, and S.~R.~Sharpe: Phys.\ Rev.\ D 38, 1278 (1988) 199
  \item M.~Creutz: Phys.\ Rev.\ D 38, 1228 (1988) 199
  \item S.~Gottlieb et al.: Phys.\ Rev.\ D 35, 2531 (1987) 199, 200
  \item G.~G.~Batrouni et al.: Phys.\ Rev.\ D 32, 2736 (1985) 199
  \item A.~Hasenfratz and A.~Alexandru: Phys.\ Rev.\ D 65, 114506 (2002) 200
  \item A.~Hasenfratz, P.~Hasenfratz, and F.~Niedermayer: Phys.\ Rev.\ D 72,
        114508 (2005) 200
  \item P.~de~Forcrand and T.~Takaishi: Nucl.\ Phys.\ B (Proc.\ Suppl.) 53,
        968 (1997) 201
  \item R.~Frezzotti and K.~Jansen: Phys.\ Lett.\ B 402, 328 (1997) 201
  \item T.~Takaishi and P.~de~Forcrand: Int.\ J.\ Mod.\ Phys.\ C 13, 343
        (2002) 201
  \item I.~Horv\'ath, A.~D.~Kennedy, and S.~Sint: Nucl.\ Phys.\ B (Proc.\
        Suppl.) 73, 834 (1999) 201
  \item M.~A.~Clark and A.~D.~Kennedy: Nucl.\ Phys.\ B (Proc.\ Suppl.) 129,
        850 (2004) 201
  \item M.~A.~Clark, P.~de~Forcrand, and A.~D.~Kennedy: PoS LAT2005, 115
        (2005) 201
  \item M.~A.~Clark: PoS LAT2006, 004 (2006) 201
  \item M.~Hasenbusch: Phys.\ Rev.\ D 59, 054505 (1999) 201
  \item M.~Hasenbusch: Phys.\ Lett.\ B 519, 177 (2001) 202
  \item P.~de~Forcrand: Nucl.\ Phys.\ B (Proc.\ Suppl.) 73, 822 (1999) 202
  \item M.~L\"uscher: Nucl.\ Phys.\ B 418, 637 (1994) 202
  \item M.~L\"uscher: JHEP 0305, 052 (2003) 203
  \item M.~L\"uscher: Comput.\ Phys.\ Commun.\ 165, 199 (2005) 203
  \item M.~L\"uscher: Comput.\ Phys.\ Commun.\ 156, 209 (2004) 203
  \item A.~Ukawa: Nucl.\ Phys.\ B (Proc.\ Suppl.) 106, 195 (2002) 203
  \item M.~Hasenbusch: Nucl.\ Phys.\ B (Proc.\ Suppl.) 129--130, 27 (2004) 203
  \item G.~Bhanot, U.~M.~Heller, and I.~O.~Stamatescu: Phys.\ Lett.\ B 129,
        440 (1983) 204
  \item M.~Grady: Phys.\ Rev.\ D 32, 1496 (1985) 204
  \item M.~Creutz: \emph{Quantum Fields on the Computer} (World Scientific,
        Singapore 1992) 204
  \item F.~Fucito, E.~Marinari, G.~Parisi, and C.~Rebbi: Nucl.\ Phys.\ B 180,
        369 (1981) 205
  \item K.~G.~Wilson: Phys.\ Rev.\ D 10, 2445 (1974) 207
  \item K.~G.~Wilson: in \emph{New Phenomena in Subnuclear Physics}, edited by
        A.~Zichichi (Plenum, New York 1977) 207
  \item S.~D\"urr et al.: Science 322, 1224 (2008) 209, 210
\end{enumerate}
